\documentclass[aspectratio=169,10pt]{beamer}
\usepackage[style=ieee,backend=biber]{biblatex}
\addbibresource{reference.bib}
\usepackage{colortbl,tabularx,mathrsfs,calligra}
\usepackage{amsmath,amsfonts,amssymb,amsthm}
\usepackage{cancel}
\usepackage{ragged2e}
\usepackage[indonesian]{babel}
\usepackage{tikz}
\usepackage{caption}
\usepackage{wrapfig}
\usepackage{multirow}
\usepackage{multicol}
\usepackage{array}
\usepackage{pgfplots, tkz-euclide,calc}
\pgfplotsset{compat=1.18}
\usepackage{listings}
\usepackage{ulem}
\usepackage{hyperref}
\usepackage{arev}
% \usepackage{pifont}
% \usepackage[scaled]{uarial}
% \renewcommand*\familydefault{\sfdefault}
\usepackage[T1]{fontenc}

\graphicspath{{C:/Users/teoso/OneDrive/Documents/Tugas Kuliah/Template Math Depart/}{./foto/}}

\definecolor{HIMAmuda}{HTML}{01D1FD}
\definecolor{HIMAtua}{HTML}{02016A}
\definecolor{HIMAabu}{HTML}{CBCBCC}
\definecolor{PastelGreen}{HTML}{77DD77}

\usetheme{Madrid}

\setbeamercolor{palette primary}{bg=HIMAtua,fg=white}
\setbeamercolor{palette secondary}{bg=HIMAmuda,fg=black}
\setbeamercolor{palette tertiary}{bg=HIMAabu,fg=black}
\setbeamercolor{palette quaternary}{bg=HIMAmuda,fg=white}
\setbeamercolor{structure}{fg=HIMAmuda} % itemize, enumerate, etc
\setbeamercolor{section in toc}{fg=HIMAtua} % TOC sections
\setbeamercolor{bibliography item}{parent=palette secondary}
\setbeamercolor*{bibliography entry author}{parent=section in toc}

\usetikzlibrary{shapes.geometric, arrows, calc, intersections}

\tikzstyle{startstop} = [ellipse, minimum width=1cm, minimum height=1cm,text centered, draw=black, fill=red!30]
\tikzstyle{process} = [rectangle, minimum width=2cm, minimum height=1cm, text centered, draw=black, fill=blue!30]
\tikzstyle{decision} = [diamond, minimum width=1cm, minimum height=1cm, text centered, draw=black, fill=blue!50]
\tikzstyle{arrow} = [thick,->,>=stealth]

\newcolumntype{L}[1]{>{\raggedright\let\newline\\\arraybackslash\hspace{0pt}}m{#1}}
\newcolumntype{C}[1]{>{\centering\let\newline\\\arraybackslash\hspace{0pt}}m{#1}}
\newcolumntype{R}[1]{>{\raggedleft\let\newline\\\arraybackslash\hspace{0pt}}m{#1}}

% \usefonttheme{professionalfonts}
\setbeamertemplate{theorems}[numbered]
\setbeamertemplate{bibliography item}[book]

\theoremstyle{definition}
\newtheorem{definisi}{Definisi}
\newtheorem{proposisi}{Proposisi}
\newtheorem{teorema}{Teorema}
\newtheorem{lema}{Lema}
\newtheorem{akibat}{Akibat}
\newtheorem{contoh}{Contoh}
\newtheorem{latihan}{Latihan}

\usepackage[nameinlink]{cleveref}
\crefname{teorema}{Teorema}{Teorema}
\crefname{proposisi}{Proposisi}{Proposisi}
\crefname{definisi}{Definisi}{Definisi}
\crefname{lema}{Lema}{Lema}
\crefname{akibat}{Akibat}{Akibat}
\crefname{contoh}{Contoh}{Contoh}
\crefname{latihan}{Latihan}{Latihan}

\newcommand{\myref}[1]{\hyperref[#1]{\nameCref{#1} \ref*{#1}}}

\newcommand{\R}{\mathbb{R}}
\newcommand{\N}{\mathbb{N}}
\newcommand{\Z}{\mathbb{Z}}
\newcommand{\C}{\mathbb{C}}

\AtBeginEnvironment{contoh}{%
\setbeamercolor{block title}{use=example text,fg=white,bg=example text.fg!75!black}
\setbeamercolor{block body}{parent=normal text,use=block title example,bg=block title example.bg!10!bg}
\setbeamercolor{item}{fg=example text.fg}
}
\AtBeginEnvironment{definisi}{
\setbeamercolor{block title}{fg=white,bg=HIMAtua}
\setbeamercolor{block body}{parent=normal text,bg=HIMAtua!30!white}
\setbeamercolor{item}{fg=HIMAtua}
}
\AtBeginEnvironment{latihan}{%
  \setbeamercolor{block title}{fg=white,bg=yellow!30!orange}
  \setbeamercolor{block body}{parent=normal text,bg=yellow!30!white}
  \setbeamercolor{item}{fg=yellow!30!orange}
}
\AtBeginEnvironment{teorema}{
    \setbeamercolor{block title}{bg=darkgray,fg=white}
    \setbeamercolor{block body}{parent=pallette tertiary,bg=HIMAabu!30!white}
    \setbeamercolor{item}{fg=HIMAabu}
}
\AtBeginEnvironment{proposisi}{
    \setbeamercolor{block title}{bg=teal,fg=white}
    \setbeamercolor{block body}{parent=pallette tertiary,bg=teal!30!white}
    \setbeamercolor{item}{fg=teal}
}
\AtBeginEnvironment{lema}{
    \setbeamercolor{block title}{bg=lime!70!black,fg=white}
    \setbeamercolor{block body}{parent=pallette tertiary,bg=lime!30!white}
    \setbeamercolor{item}{fg=lime}
}
\AtBeginEnvironment{akibat}{
    \setbeamercolor{block title}{bg=magenta,fg=white}
    \setbeamercolor{block body}{parent=pallette tertiary,bg=magenta!30!white}
    \setbeamercolor{item}{fg=magenta}
}

\date{25 Mei 2026}
\title[Kalkulus Itô]{Kalkulus Stokastik Itô: Definisi dan Contoh}
\author[Tetew]{Teosofi Hidayah Agung}
\institute[Matematika ITS]{Matematika ITS}


\begin{document}


\begin{frame}
  \titlepage
\end{frame}

\AtBeginSection{
  {
      \begin{frame}{Daftar Isi}
        \tableofcontents[currentsection]
      \end{frame}}
}

\section{Kalkulus Itô}

\begin{frame}{Apa itu Kalkulus Itô?}
  \begin{itemize}
    \item \textbf{Kalkulus Itô} adalah perluasan dari kalkulus klasik untuk fungsi-fungsi stokastik, khususnya proses yang memiliki lintasan kasar (tidak terdiferensiasi di sembarang titik) seperti Gerak Brown (Proses Wiener).
    \item \textbf{Penemu:} Diperkenalkan dan dikembangkan secara bertahap oleh matematikawan Jepang \textbf{Kiyosi Itô} pada pertengahan abad ke-20.
    \item \textbf{Kegunaan:}
          \begin{itemize}
            \item Matematika Keuangan: Model harga opsi (Model Black-Scholes), pemodelan suku bunga.
            \item Fisika: Gerak partikel acak, mekanika kuantum.
            \item Biologi dan rekayasa: Dinamika populasi acak dan pengendalian stokastik.
          \end{itemize}
  \end{itemize}
\end{frame}

\section{Integral Stokastik Itô}

\begin{frame}{Latar Belakang Integral Stokastik}
  Penjumlahan Riemann yang digunakan untuk mendekati integral suatu fungsi $f : [0, T] \to \R$ ditulis sebagai:
  \begin{equation*}
    \sum_{j=0}^{n-1} f(s_j)(t_{j+1} - t_j),
  \end{equation*}
  dengan $0 = t_0 < t_1 < \dots < t_n = T$ dan $s_j \in [t_j, t_{j+1}]$. Nilai integral klasik tidak bergantung pada pilihan $s_j$.
  \vspace{0.3cm}

  Namun, dalam kasus stokastik, bentuk penjumlahan pendekatannya adalah:
  \begin{equation*}
    \sum_{j=0}^{n-1} f(s_j)(W(t_{j+1}) - W(t_j)),
  \end{equation*}
  dengan $W(t)$ adalah proses Wiener. Pada integral stokastik limit jumlah tersebut \textbf{sangat bergantung} pada pilihan titik $s_j$. Agar nilainya tidak bergantung pada kejadian di masa depan, kita pilih $s_j = t_j$.
\end{frame}

\begin{frame}{Proses Sederhana Acak}
  \begin{definisi}[5.1]
    Suatu fungsi stokastik $f(t)$, dengan $t \geq 0$, disebut sebagai \textbf{proses sederhana acak} jika terdapat suatu partisi berhingga
    \begin{equation*}
      0 = t_0 < t_1 < \dots < t_n,
    \end{equation*}
    dan variabel acak $\eta_0, \eta_1, \dots, \eta_{n-1}$ terintegralkan kuadrat, yaitu $\mathbb{E}[|\eta_j|^2] < \infty$, sehingga untuk setiap $t \geq 0$ berlaku
    \begin{equation*}
      f(t) = \sum_{j=0}^{n-1} \eta_j \mathbb{1}_{[t_j, t_{j+1})}(t),
    \end{equation*}
    dengan setiap $\eta_j$ merupakan variabel acak yang $\mathcal{F}_{t_j}$-terukur.
    Himpunan semua proses sederhana acak tersebut dilambangkan dengan $M_{\text{step}}^2$.
  \end{definisi}
\end{frame}

\begin{frame}{Definisi Integral Stokastik Itô}
  \begin{definisi}[5.2]
    Integral stokastik dari suatu proses sederhana acak $f \in M_{\text{step}}^2$, yang memiliki representasi
    \begin{equation*}
      f(t) = \sum_{j=0}^{n-1} \eta_j \mathbb{1}_{[t_j, t_{j+1})}(t),
    \end{equation*}
    dengan $\eta_j$ adalah variabel acak yang $\mathcal{F}_{t_j}$-terukur dan  $\mathbb{E}[|\eta_j|^2] < \infty$, didefinisikan sebagai
    \begin{equation*}
      I^*(f) = \sum_{j=0}^{n-1} \eta_j (W(t_{j+1}) - W(t_j)).
    \end{equation*}
  \end{definisi}
\end{frame}

\begin{frame}{Sifat Integral Stokastik}
  \begin{proposisi}[5.1]
    Untuk setiap proses sederhana acak $f \in M_{\text{step}}^2$, integral stokastik $I^*(f)$ merupakan variabel acak yang dapat diintegralkan secara kuadrat, yaitu $I^*(f) \in L^2$ dan berlaku:
    \begin{equation*}
      \mathbb{E}((I^*(f))^2) = \mathbb{E}\left(\int_0^\infty f(t)^2 \, dt\right).
    \end{equation*}
  \end{proposisi}
\end{frame}

\begin{frame}[allowframebreaks]{Sifat Ekspektasi Bersyarat}
  \begin{proposisi}[5.2]
    Ekspektasi bersyarat memiliki beberapa sifat berikut:
    \begin{enumerate}
      \item Ekspektasi bersyarat bersifat linear, yaitu $\mathbb{E}(\alpha \xi + \beta \zeta | \mathcal{G}) = \alpha \mathbb{E}(\xi | \mathcal{G}) + \beta \mathbb{E}(\zeta | \mathcal{G})$;
      \item Ekspektasi dari nilai ekspektasi bersyarat sama dengan ekspektasi awal, yaitu $\mathbb{E}(\mathbb{E}(\xi | \mathcal{G})) = \mathbb{E}(\xi)$;
      \item Jika $\xi$ dapat diukur terhadap $\mathcal{G}$, maka ekspektasi bersyaratnya sama dengan $\xi$, yaitu $\mathbb{E}(\xi \zeta | \mathcal{G}) = \xi \mathbb{E}(\zeta | \mathcal{G})$;
      \item Jika $\xi$ independen terhadap $\mathcal{G}$, maka $\mathbb{E}(\xi | \mathcal{G}) = \mathbb{E}(\xi)$;
      \item Memenuhi sifat menara, yaitu $\mathbb{E}(\mathbb{E}(\xi | \mathcal{G}) | \mathcal{H}) = \mathbb{E}(\xi | \mathcal{H})$ jika $\mathcal{H} \subset \mathcal{G}$;
      \item Jika $\xi \geq 0$, maka $\mathbb{E}(\xi | \mathcal{G}) \geq 0$.
    \end{enumerate}
  \end{proposisi}
\end{frame}

\begin{frame}[allowframebreaks]{Bukti Proposisi 5.1}
  Misalkan inkremen $\Delta_j W = W(t_{j+1}) - W(t_j)$ dan $\Delta_j t = t_{j+1} - t_j$. Maka $\mathbb{E}(\Delta_j W) = 0$ dan $\mathbb{E}((\Delta_j W)^2) = \Delta_j t$.\\
  Pertama dihitung ekspektasinya:
  \begin{equation*}
    (I^*(f))^2 = \sum_{j=0}^{n-1}\sum_{k=0}^{n-1} \eta_j \eta_k \Delta_j W \Delta_k W.
  \end{equation*}
  Sehingga didapatkan ekspektasinya menjadi:
  \begin{equation*}
    \mathbb{E}((I^*(f))^2) = \sum_{j=0}^{n-1} \mathbb{E}(\eta_j^2 (\Delta_j W)^2) + 2 \sum_{k < j} \mathbb{E}(\eta_j \eta_k \Delta_j W \Delta_k W).
  \end{equation*}

  Karena $\Delta_j W$ independen terhadap $\mathcal{F}_{t_j}$ dan $\eta_j$ terukur pada $\mathcal{F}_{t_j}$, berlaku:
  \begin{align*}
    \mathbb{E}(\eta_j^2 (\Delta_j W)^2) & = \mathbb{E}(\mathbb{E}(\eta_j^2 (\Delta_j W)^2 | \mathcal{F}_{t_j}))                \\
                                        & = \mathbb{E}(\eta_j^2 \mathbb{E}((\Delta_j W)^2)) = \mathbb{E}(\eta_j^2) \Delta_j t.
  \end{align*}

  Demikian juga untuk suku silang dengan $k < j$:
  \begin{align*}
    \mathbb{E}(\eta_j \eta_k \Delta_j W \Delta_k W) & = \mathbb{E}(\mathbb{E}(\eta_j \eta_k \Delta_j W \Delta_k W | \mathcal{F}_{t_j})) \\
                                                    & = \mathbb{E}(\eta_j \eta_k \Delta_k W \mathbb{E}(\Delta_j W | \mathcal{F}_{t_j})) \\
                                                    & = 0.
  \end{align*}

  Karena suku silangnya nol, kita simpulkan:
  \begin{equation*}
    \mathbb{E}((I^*(f))^2) = \sum_{j=0}^{n-1} \mathbb{E}(\eta_j^2) \Delta_j t.
  \end{equation*}
  Karena $\eta_j \in L^2$, maka $I^*(f) \in L^2$.

  Di sisi lain,
  \begin{equation*}
    \int_0^\infty f(t)^2 \, dt = \sum_{j=0}^{n-1} \eta_j^2 \mathbb{1}_{[t_j, t_{j+1})}(t).
  \end{equation*}
  Sehingga:
  \begin{equation*}
    \mathbb{E}\left(\int_0^\infty f(t)^2 \, dt\right) = \sum_{j=0}^{n-1} \mathbb{E}(\eta_j^2) \Delta_j t.
  \end{equation*}
  Kedua sisi memiliki nilai yang sama, sehingga berlaku $\mathbb{E}((I^*(f))^2) = \mathbb{E}\left(\int_0^\infty f(t)^2 \, dt\right)$. \qed
\end{frame}

\begin{frame}{Lema 5.1}
  \begin{lema}[5.1]
    Untuk setiap proses sederhana acak $f, g \in M_{\text{step}}^2$, berlaku
    \begin{equation*}
      \mathbb{E}(I^*(f)I^*(g)) = \mathbb{E}\left(\int_0^\infty f(t)g(t) \, dt\right).
    \end{equation*}
  \end{lema}
\end{frame}

\begin{frame}[allowframebreaks]{Bukti Lema 5.1}
  Misalkan $f, g \in M_{\text{step}}^2$. Dengan demikian terdapat suatu partisi berhingga yang sama (setelah penghalusan partisi),
  \begin{equation*}
    f(t) = \sum_{j=0}^{n-1} \eta_j \mathbb{1}_{[t_j, t_{j+1})}(t), \quad g(t) = \sum_{j=0}^{n-1} \zeta_j \mathbb{1}_{[t_j, t_{j+1})}(t).
  \end{equation*}
  Kemudian,
  \begin{equation*}
    \mathbb{E}(I^*(f)I^*(g)) = \sum_{j=0}^{n-1}\sum_{k=0}^{n-1} \mathbb{E}(\eta_j \zeta_k \Delta_j W \Delta_k W).
  \end{equation*}

  Untuk kasus $j = k$:
  \begin{align*}
    \mathbb{E}(\eta_j \zeta_j (\Delta_j W)^2) & = \mathbb{E}(\eta_j \zeta_j \mathbb{E}((\Delta_j W)^2 | \mathcal{F}_{t_j})) \\
                                              & = \mathbb{E}(\eta_j \zeta_j) \Delta_j t.
  \end{align*}

  Untuk kasus $j \neq k$, asumsikan $j < k$:
  \begin{align*}
    \mathbb{E}(\eta_j \zeta_k \Delta_j W \Delta_k W) & = \mathbb{E}(\eta_j \zeta_k \Delta_j W \mathbb{E}(\Delta_k W | \mathcal{F}_{t_k})) \\
                                                     & = 0.
  \end{align*}
  Sehingga semua suku silang nol.

  Maka tersisa $\mathbb{E}(I^*(f)I^*(g)) = \sum_{j=0}^{n-1} \mathbb{E}(\eta_j \zeta_j) \Delta_j t$.
  Sedangkan untuk ruas kanan,
  \begin{align*}
    \mathbb{E}\left(\int_0^\infty f(t)g(t) \, dt\right) & = \mathbb{E}\left( \int_0^\infty \sum_{j=0}^{n-1} \eta_j \zeta_j \mathbb{1}_{[t_j, t_{j+1})}(t) \, dt \right) \\
                                                        & = \sum_{j=0}^{n-1} \mathbb{E}(\eta_j \zeta_j) \Delta_j t.
  \end{align*}
  Kedua ruas bernilai sama. Bukti selesai. \qed
\end{frame}

\begin{frame}{Linieritas Pemetaan $I^*$}
  \begin{lema}[5.2]
    Pemetaan $I^* : M_{\text{step}}^2 \to L^2$ adalah pemetaan linear, yaitu untuk setiap $f, g \in M_{\text{step}}^2$ dan konstanta $\alpha, \beta \in \R$ berlaku:
    \begin{equation*}
      I^*(\alpha f + \beta g) = \alpha I^*(f) + \beta I^*(g).
    \end{equation*}
  \end{lema}
\end{frame}

\begin{frame}[allowframebreaks]{Bukti Lema 5.2}
  Misalkan $f, g \in M_{\text{step}}^2$, dan representasi mereka pada partisi bersama adalah
  \begin{equation*}
    f(t) = \sum_{j=0}^{n-1} \eta_j \mathbb{1}_{[t_j, t_{j+1})}(t), \quad g(t) = \sum_{j=0}^{n-1} \zeta_j \mathbb{1}_{[t_j, t_{j+1})}(t).
  \end{equation*}
  Maka untuk konstanta $\alpha, \beta$:
  \begin{align*}
    (\alpha f + \beta g)(t) & = \alpha f(t) + \beta g(t)                                                         \\
                            & = \sum_{j=0}^{n-1} (\alpha \eta_j + \beta \zeta_j) \mathbb{1}_{[t_j, t_{j+1})}(t).
  \end{align*}
  Berdasarkan definisi integral stokastik:
  \begin{align*}
    I^*(\alpha f + \beta g) & = \sum_{j=0}^{n-1} (\alpha \eta_j + \beta \zeta_j)(W(t_{j+1}) - W(t_j))                                       \\
                            & = \alpha \sum_{j=0}^{n-1} \eta_j (W(t_{j+1}) - W(t_j)) + \beta \sum_{j=0}^{n-1} \zeta_j (W(t_{j+1}) - W(t_j)) \\
                            & = \alpha I^*(f) + \beta I^*(g).
  \end{align*}
  Bukti telah ditunjukkan. \qed
\end{frame}

\section{Contoh-Contoh}

\end{document}